% !PRINTED 264
\weberpage{264}
%<<CONT
have; that, however, in the chain of the principal subdeterminants
$R'_k$ of the form $\varphi'$ no vanishing terms occur; and that,
finally, to a non-vanishing $R_k$ there corresponds an $R'_k$ of the
same sign. Then the numbers $\pi$, $\nu$ for
$\varphi$ and for $\varphi'$ can be determined both from the series (4) and from the series
(5), and the number of \gesperrt{changes of sign}, which
is the same in both, gives the number $\nu$ of the negative squares.

Now, in order to prove that the coefficient $\varepsilon$ can
be determined in the manner indicated, we assume
that $\varphi$ has somehow been transformed into a sum of squares
\[
\varphi = \lambda_1 y_1^2 + \lambda_2 y_2^2 + \cdots + \lambda_m y_m^2,
\]
and that $\psi$, expressed in the variables $y_1, y_2 \ldots y_m$, has the
expression
\[
\psi = \Sigma \beta_{i,k} y_i y_k.
\]
The numbers $\pi$, $\nu$ for the function $\varphi'$ are then determined from the
chain of the principal subdeterminants of
\[
\begin{vmatrix}
\lambda_1 + \varepsilon\beta_{1,1}, & \varepsilon\beta_{1,2} \ldots & \varepsilon\beta_{1,m} \\
\varepsilon\beta_{2,1}, & \lambda_2 + \varepsilon\beta_{2,2} \ldots & \varepsilon\beta_{2,m} \\
\cdot & \cdot & \cdot \\
\varepsilon\beta_{m,1}, & \varepsilon\beta_{m,2} \ldots & \lambda_m + \varepsilon\beta_{m,m}
\end{vmatrix}
\]
according to Theorem V.

Now, however, one may take $\varepsilon$ so small that these principal
subdeterminants agree \gesperrt{in sign} with
\[
\lambda_1\lambda_2\lambda_3 \ldots \lambda_m, \quad \lambda_1\lambda_2\lambda_3 \ldots \lambda_{m-1}, \quad \ldots, \quad \lambda_1\lambda_2, \quad \lambda_1, \quad 1,
\]
and then the number of positive and negative squares
of $\varphi'$ is equal to the number of positive and negative ones among the
coefficients $\lambda$, of which none vanishes, i.\,e. the numbers
$\pi$ and $\nu$ are the same for $\varphi$ and $\varphi'$.

If now the coefficients of $\psi$, expressed in the original
variables, are $b_{i,k}$, so that
\[
\psi = \Sigma b_{i,k} x_i x_k,
\]
then one of the principal subdeterminants of $\varphi'$ is
\[
R'_k = \begin{vmatrix}
a_{1,1} + \varepsilon b_{1,1} & \ldots & a_{1,k} + \varepsilon b_{1,k} \\
\cdot & \cdot & \cdot \\
a_{k,1} + \varepsilon b_{k,1} & \ldots & a_{k,k} + \varepsilon b_{k,k}
\end{vmatrix},
\]
and is thus an integral rational function of degree $k$ in $\varepsilon$,
\[
R'_k = R_k + \varepsilon M_1 + \varepsilon^2 M_2 + \ldots \varepsilon^k M_k,
\]
%CONT>>
